% >>> PROPAGATION (Fable 5, Aug 29 2026): one paragraph added at \S discussion placing Decomposition Neutrality beside the union
% (alternatives vs conjunctive composites; separate ledger entries in the companions). See the ledger, R109.
\documentclass[11pt]{article}
\usepackage[margin=1.1in]{geometry}
\usepackage{amsmath, amssymb, amsthm}
\usepackage[backend=biber, style=authoryear]{biblatex}
\usepackage[hidelinks]{hyperref}
\addbibresource{refs.bib}
\addbibresource{excl_refs.bib}
\newtheorem{proposition}{Proposition}
\newtheorem{theorem}{Theorem}
\newtheorem{definition}{Definition}
\newtheorem{remark}{Remark}

\usepackage{filecontents}
\usepackage{booktabs}
\begin{filecontents*}{excl_refs.bib}
@book{wittgenstein1921,
  author = {Wittgenstein, Ludwig},
  title = {Tractatus Logico-Philosophicus},
  year = {1921},
  note = {Trans. Pears and McGuinness, Routledge, 1961}
}
@article{wittgenstein1929,
  author = {Wittgenstein, Ludwig},
  title = {Some Remarks on Logical Form},
  journal = {Proceedings of the Aristotelian Society, Supplementary Volumes},
  volume = {9}, pages = {162--171}, year = {1929}
}
@article{ramsey1923,
  author = {Ramsey, Frank P.},
  title = {Critical Notice: Tractatus Logico-Philosophicus},
  journal = {Mind}, volume = {32}, pages = {465--478}, year = {1923}
}
@book{jevons1864,
  author = {Jevons, William Stanley},
  title = {Pure Logic, or the Logic of Quality apart from Quantity},
  publisher = {Edward Stanford}, year = {1864}
}
% VERIFY the six entries below against the actual publications (details per referee recall).
% NOTE on putnam1968 specifically: the colloquium was 1968 but Boston Studies vol. 5 appeared 1969;
% citations split between the two years --- check against the volume itself.
@article{birkhoffvonneumann1936,
  author = {Birkhoff, Garrett and von Neumann, John},
  title = {The Logic of Quantum Mechanics},
  journal = {Annals of Mathematics}, volume = {37}, pages = {823--843}, year = {1936}
}
@incollection{putnam1968,
  author = {Putnam, Hilary},
  title = {Is Logic Empirical?},
  booktitle = {Boston Studies in the Philosophy of Science},
  year = {1968},
  note = {Reprinted as ``The Logic of Quantum Mechanics''}
}
@article{barneswilliams2011,
  author = {Barnes, Elizabeth and Williams, J. Robert G.},
  title = {A Theory of Metaphysical Indeterminacy},
  journal = {Oxford Studies in Metaphysics}, volume = {6}, year = {2011}
}
@article{darbypickup2021,
  author  = {Darby, George and Pickup, Martin},
  title   = {Modelling deep indeterminacy},
  journal = {Synthese},
  volume  = {198}, pages = {1685--1710}, year = {2021}
}
@article{calosiwilson2021,
  author  = {Calosi, Claudio and Wilson, Jessica},
  title   = {Quantum indeterminacy and the double-slit experiment},
  journal = {Philosophical Studies},
  volume  = {178}, year = {2021},
  note    = {verify title, volume and pages at merge}
}
@article{darby2010,
  author = {Darby, George},
  title = {Quantum Mechanics and Metaphysical Indeterminacy},
  journal = {Australasian Journal of Philosophy}, volume = {88}, year = {2010}
}
@article{skow2010,
  author = {Skow, Bradford},
  title = {Deep Metaphysical Indeterminacy},
  journal = {The Philosophical Quarterly}, volume = {60}, year = {2010}
}
@article{funkhouser2006,
  author = {Funkhouser, Eric},
  title = {The Determinable--Determinate Relation},
  journal = {No\^us}, volume = {40}, year = {2006}
}
@book{boole1854,
  author = {Boole, George},
  title = {An Investigation of the Laws of Thought},
  year = {1854}, publisher = {Walton and Maberly}
}
@article{wilson2013,
  author = {Wilson, Jessica M.},
  title = {A Determinable-Based Account of Metaphysical Indeterminacy},
  journal = {Inquiry}, volume = {56}, pages = {359--385}, year = {2013}
}
@article{calosi2019,
  author = {Calosi, Claudio and Wilson, Jessica M.},
  title = {Quantum Metaphysical Indeterminacy},
  journal = {Philosophical Studies}, volume = {176}, pages = {2599--2627}, year = {2019}
}
@misc{Aethus, howpublished={Manuscript}, author={Benander, Ajax}, title={Aethic Reasoning: A Comprehensive Solution to the Quantum Measurement Problem}, year={2024}, month = {11}, url={https://philpapers.org/rec/BENARA-5}, eprint={https://philsci-archive.pitt.edu/id/eprint/25713}, eprinttype={Alternative URL}, note={Manuscript; the complete development of the framework, successive deposits November 2024--2026}}
@misc{AethusQuant, author={Benander, Ajax}, title={Quantum Mechanics from Records: Measurement Structure on a Finite Path-Kernel Domain}, year={2026}, howpublished={Preprint}, url={https://philpapers.org/rec/BENTQB}}
@misc{AethusBrief, author={Benander, Ajax}, title={Aethic Reasoning in Brief: From a Counterfactual Semantics to the Quantum Measurement Problem}, year={2026}, howpublished={Preprint}, url={https://philpapers.org/rec/BENARI-3}, eprint={https://philsci-archive.pitt.edu/id/eprint/30439}, eprinttype={Alternative URL}}
@article{fine2017truthmaker,
  author = {Fine, Kit},
  title = {A theory of truthmaker content {I}: Conjunction, disjunction and negation},
  journal = {Journal of Philosophical Logic},
  volume = {46}, number = {6}, pages = {625--674}, year = {2017}
}
@article{humberstone1981from,
  author = {Humberstone, I. L.},
  title = {From worlds to possibilities},
  journal = {Journal of Philosophical Logic},
  volume = {10}, number = {3}, pages = {313--339}, year = {1981}
}
@incollection{holliday2021possibility,
  author = {Holliday, Wesley H.},
  title = {Possibility semantics},
  booktitle = {Selected Topics from Contemporary Logics},
  editor = {Fitting, Melvin},
  publisher = {College Publications}, address = {London}, year = {2021}
}
@misc{AethusTree,
  author       = {Benander, Ajax},
  title        = {The Weight Algebra of Aethae: A Two-Layer Commutative Semiring for Superposition and Pruning},
  year         = {2026},
  howpublished = {Preprint},
  url={https://philpapers.org/rec/BENTWA-5},
  note         = {The algebraic support extraction; owns the Aethus as a formal
                  object, the refinement order, formal validity, and the two-layer algebra.}
}

@misc{GoldenScience2022,
  author       = {Benander, Ajax},
  title        = {Golden Science},
  year         = {2022},
  month        = {November},
  howpublished = {Self-hosted manuscript},
  url          = {https://aethus.org/aethic-code/files/GoldenScience_Nov2022.pdf},
  note         = {Title page August 2022, the date of the first version of the source; this version created November 2022. Not deposited to a
                  preprint server; cited as a dated developmental record rather than
                  as a source of results.}
}
\end{filecontents*}

\title{Exclusivity as Achievement:\\ \large Relocating Mutual Exclusion to the Interior, and the Composition Regime That Falls Out}
\author{Ajax Benander\thanks{Incoming graduate student, Department of Physics, Brown University (Fall 2026). Email: \href{mailto:akb9408@rit.edu}{akb9408@rit.edu}, Website: \url{https://aethus.org}. The framework drawn on here is the author's, developed 2021--2026 and documented in the companion works \cite{AethusBrief, Aethus}; AI assistance in drafting the present manuscript to conventional standards is disclosed per venue policy.}}
\date{First drafted July 2026\\ {\small This version: September 2026 --- draft}}

\begin{document}
\maketitle

\begin{abstract}
This paper argues that the mutual exclusivity of alternatives is structural content of the alternatives themselves rather than an intrinsic feature of disjunction. Granting a primitive composition of alternatives, the composition anyone who takes alternatives to compose already attributes to them, we encode exclusion in the interior content of the alternatives and classify the resulting exclusion profiles: the fully exclusive regime, in which co-present alternatives carry mutual negations, and the agreeing regime, in which they do not, the latter obtained from the same composition without any regime-specific primitive. A one-statement test establishes that the two forms are nonidentical within the structured representation, and profiles with different exclusion content are shown to update differently under the same operations. We distinguish exclusion profiles from the irreducibility marker that agreeing superposition in its irreducible form requires, and both from the question whether the world instantiates the agreeing regime, which is a separate physical proposal on which no thesis here depends. We argue that the proposal reclassifies the color-exclusion problem of the \emph{Tractatus} and offers the determinable-based literature on metaphysical indeterminacy a composition story distinct from fusion and restricted persistence.
\end{abstract}

\section{The Question, and Its Century of Standing}\label{sec:intro}
Take the oldest small embarrassment in analytic philosophy: \emph{this is red} and \emph{this is green}, said of one point at one time, exclude each other --- yet the exclusion is neither a logical contradiction, the sentences sharing no negated conjunct, nor comfortably empirical, since no experience seems needed to know it. Wittgenstein legislated the difficulty away in the Tractatus (6.3751) by promising that a completed analysis would reveal the incompatibility as tautologous \parencite{wittgenstein1921}; Ramsey pressed the promise \parencite{ramsey1923}; the attempted repair \parencite{wittgenstein1929} conceded that elementary propositions can exclude one another without contradicting one another --- and the concession is standardly credited with collapsing the Tractarian program, since a logic in which atoms exclude is no longer doing its exclusions \emph{by logic}. The repair literature that followed --- from the 1929 paper's own picture of coordinated scales to later enriched-form and property-structure treatments --- kept relocating the exclusion \emph{within} logical form or the metaphysics of properties; the present move differs in kind, relocating it into the alternatives' stated content.

The century since has largely filed the episode as a lesson about logical form: analysis is harder than the Tractatus hoped, atomism was too optimistic, exclusion must be built into an enriched language or accepted as metaphysical necessity or handed to the determinate--determinable structure of properties \parencite{funkhouser2006}.

\textit{What every response took for granted.} Each response takes for granted the thing this paper puts in question --- that exclusion, wherever we lodge it, is a fixed feature of the alternatives' \emph{nature}, something they could not have lacked. And the stakes are larger than one puzzle in the philosophy of logic: on the view to be developed, the century of color-exclusion and the concurrent century of quantum-foundations bewilderment are the \emph{same} confusion seen from two sides --- exclusivity mistaken for the nature of \emph{or} --- so that a single relocation addresses both, (the second only gestured at here --- the companions carry it, with Section~\ref{sec:exposure} documenting the provenance that links the two). The proposal here is that the datum resisted classification for a principled reason: exclusion is neither logical form nor brute matter because it is \emph{stated structure} --- content the alternatives carry, an achievement rather than a nature, and therefore, in the strict sense that will matter throughout, \emph{removable} in the representation. What the paper establishes directly is a representation in which exclusionary structure can vary; whether actual properties admit that variation is the further argument of \S\ref{sec:color}, conditional on how properties are individuated. Each of the standard responses then turns out to hold half of the view: the enriched-language response is right that exclusion is contentful, wrong that the content is semantic rather than worldly; the necessitarian is right that exclusion, where present, is structural rather than accidental, wrong that structure entails ineliminability; and the determinable theorist is right about the shape of the states involved, while lacking --- as Section~\ref{sec:taxon} argues --- an account of how such states \emph{compose}.



Relative to the framework overview \cite{AethusBrief}, whose union section compresses the present material to two pages, this standalone's delta is threefold and deliberate: the color-exclusion problem treated as a full application rather than an aside; the Jevons--Boole history developed at depth, as the reform's precedent rather than its ornament; and direct engagement with the determinable-based indeterminacy literature, which the overview only gestures at. A reader arriving from the overview should read for those three.

\section{The Relocation Move}\label{sec:relocation}
Write co-present alternatives $A_1, \dots, A_n$ and distinguish two compositional forms:
\begin{equation}\label{eq:forms}
\text{joint:}\;\; \bigcup_i a_i \qquad\qquad \text{exclusive:}\;\; \bigcup_i \Big( a_i \cap \bigcap_{j \neq i} \neg a_j \Big).
\end{equation}
The lower-case letters are \emph{pins}, the options of one attribute, and the alternatives are \emph{states}: in the exclusive form each term $a_i \cap \bigcap_{j \ne i} \neg a_j$ is the state $[X : \{a_i\}]$, whose interior holds the pin $a_i$ and, filling $X$ with exactly that pin, excludes the rest; the negations are exclusions of pins, the intersections that carry them are interiors, and the outer union composes states. In the joint form the pins are composed within one interior, the state $[X : \{a_1, \ldots, a_n\}]$ holding all of them. The display is motivation rather than a typed identity: its joint side is a single filling, one state holding several pins, while the results below concern structured compositions $J(A_1, \ldots, A_n)$ of alternatives whose interiors can be edited, the two constructions being connected only as far as Definition~\ref{def:irreducible-level} supplies: a carved filling corresponds to the composition of its determinate cells, and the marked filling is kept distinct from the profile results. An interior is always of a state, and what lives in it are pins, never other states or their negations; throughout the paper an \emph{interior negation} means exactly this, a pin omitted from a filling, the word kept because it names the omission's logical force.

The union sign is the composition of alternatives---the operation by which several alternatives are made contents of one composite---and not the extensional union of sets of worlds, on which $A \cup B$ would say only that an element lies in $A$ or in $B$; what it denotes is the co-presence of the composed alternatives in one object. Two constructions wear the sign, and they are kept distinct from here on, since a display that passes from one to the other changes the type of its operands. Assembling pins into \emph{one filling}, $[X : \{a, b\}]$, composes alternatives within an attribute, and is written with the set union inside the bracket; collecting several Aethae into a \emph{support}, $\{[X : \{a\}], [X : \{b\}]\}$, composes states, and is written with $+$ between them. The support semiring below is the algebra of the second construction; it does not make the first identical to it, and the expansion of a filling into the support of its carved pieces is not an equality in that semiring but the carving-induced passage between the two.

What makes the operation \emph{or}-like beyond the notation is its algebra, and the algebra is inherited rather than stipulated here: in the companion framework the composition of unweighted states is the addition of a commutative semiring---commutative, associative, and idempotent, $A \cup A = A$, with joint assertion as the multiplication that distributes over it---so that the idempotent law Jevons pressed on Boole is, on the world side, a theorem of the companion algebra and not a further posit \parencite{AethusTree}. The weighted extension of that algebra, in which alternatives carry measures and the composition sums them, drops idempotence and keeps the rest; the present paper's claims are stated at the unweighted layer, and its carrier should be named, since it is not the companion's intrinsic $U$, whose addition is the core. It is the support semiring $\mathcal{P}_{\mathrm{fin}}(U)$: finite sets of Aethae, added by union, $S + T = S \cup T$, and multiplied pointwise, $ST = \{ab : a \in S,\ b \in T\}$, an idempotent sum that keeps its terms and over which joint assertion distributes exactly. It is the image of the weighted algebra under forgetting weights over $\mathbb{R}_{\ge 0}$, each term's Aethus surviving and its weight dropped; over a coefficient semiring with zero divisors, $M = \mathbb{R}_{\ge 0} \times \mathbb{R}_{\ge 0}$ for instance, where $(1, 0)(0, 1) = 0$, the product of two nonempty supports need not be the support of the product, so the identification carries a no-zero-divisors hypothesis; invalidity is represented there by removing invalid members, and an alternative's structured interior is the Aethus it is. The agreeing regime is the sum's unobstructed case. The composition is not the order-theoretic join: the join of two states in the refinement order is their core, their last common ancestor, and joint assertion distributes over the union exactly but not over the core.

The placement of the negations is the entire content. In the exclusive form each alternative's \emph{interior} excludes its rivals' pins, by holding its own pin alone: exclusion is content \emph{of} the alternatives, part of what each one says. ``Says'' is a relation on the worldly side and it is owed a definition rather than a metaphor. An alternative is a structured object $A = \langle \mathrm{ch}(A),\ \mathrm{Int}(A) \rangle$: a positive character, and an interior, the latter what fills the attributes the alternative states: for each such attribute $X$, the set $S$ of pins with which $A$ fills it, the state $[X : S]$; in the companion's vocabulary the interior is the alternative's stated-attribute set, each stated attribute given with its filling, and the character its remaining content \parencite{Aethus}. Because $[X : S]$ fills $X$ with exactly the pins of $S$, it \emph{excludes} every pin of $X$ outside $S$. $A$ \emph{excludes} a pin $p$ iff $A$ states $p$'s attribute with a filling that omits $p$; $A$ excludes $B$ iff $A$ excludes some pin $B$ holds; and $A$ and $B$ \emph{exclude} iff each excludes a pin of the other. That is a primitive relation---a pin's membership in a filling---and the paper does not pretend otherwise; what has been relocated is exclusion, from a primitive incompatibility between alternatives into a primitive of a more productive kind, stated content that can be present or absent. On this definition Proposition~\ref{prop:onestatement} below is a theorem about structured objects and not a figure of speech.

In the joint form the interiors are silent, and nothing excludes. Classical metaphysics reads the left form as incoherent and the right form as what \emph{or} simply means; the relocation move reads the right form as the left form \emph{plus stated interior structure} --- so that exclusivity, far from being the nature of disjunction, is an addition to it.

The move has a precedent, and acknowledging it sharpens what is and is not new. Within logic proper, the question whether \emph{or} is inclusive or exclusive was settled by Jevons's reform of Boole's algebra \parencite{jevons1864}: Boole's own ``$+$'' was effectively exclusive, defined only for disjoint classes \parencite{boole1854}, and it was Jevons's law $X + X = X$ that made disjunction inclusive --- a reform Boole himself rejected when Jevons pressed it in correspondence, which is its own evidence of how deep the exclusive habit ran. On the reformed algebra, exclusivity, where wanted, is written in --- $A \vee B$ against $(A \vee B) \wedge \neg(A \wedge B)$ --- an added conjunct, never a feature of the connective. Logic absorbed the lesson so thoroughly that the exclusive reading now requires its own symbol. Metaphysics absorbed it unevenly at best: mereology, truthmaker and situation semantics, and the indeterminacy literature have each investigated overlapping, fused, or non-exclusive structures, and the contrast drawn here is with the treatment of \emph{alternatives} specifically, where the exclusive reading has held. When the disjuncts are not sentences in a calculus but \emph{alternatives of reality} --- ways a thing might be, outcomes a situation might take --- the exclusive reading has remained the unmarked default, treated not as one composition among two but as what alternative-hood consists in. The relocation move is Jevons's reform performed one level down: for the alternatives themselves, exclusion is an added statement, and the inclusive composition is the unmarked case.
% >>> EDIT U2 (Fable 5, Aug 29 2026): the precedent scoped --- invisibility, not warrant
The precedent's scope should be fixed before it is leaned on, since it is easy to over-read. Jevons's was a definitional choice about an algebraic operation, motivated by internal tidiness and settled by the algebra's own convenience; the relocation is a claim about alternatives of reality, and no amount of algebraic precedent establishes it. What the precedent supplies is an explanation of the reform's \emph{invisibility} --- that a composition habit this deep can go unexamined for a century and a half, in a field expressly devoted to examining such habits --- and not a warrant for performing it again one level down. The warrant is the argument of the sections that follow.
% <<< END EDIT U2
Two consequences follow immediately, and they are the paper's spine. First, a definition:
\begin{definition}[Or-jointness]\label{def:jointness}
Let composition of alternatives be primitive. Or-jointness is that composition in its \emph{unobstructed} regime: the regime in which no alternative's interior excludes a pin another's holds. Exclusivity is the same composition with such exclusions stated.
\end{definition}
\noindent Two things in the definition are distinct, because running them together is the slide the argument must not make. The composition is primitive, and the paper does not derive it from anything: it is the worldly composition attributed to alternatives of reality by anyone who takes them to compose at all, and its status is the one the companion framework gives it as a founding principle---\emph{jointness is primitive and exclusion is an achievement} \parencite{Aethus}. What is \emph{not} a new primitive is the distinction between the regimes. That distinction is fixed by logical structure alone---by which negations are and are not stated on the interior side, with no further physical, semantic, or psychological parameter anywhere in its specification---and it is in that sense deflationary. So nothing is added to the world to obtain the agreeing regime in its bare form, the absence of interior exclusion, beyond the composition one already had; what is declined is the classical content, and the irreducible form is a further entry, the marker, not a further declension. The inference the relocation licenses is therefore this. A skeptic who grants the relocation grants that exclusion is interior content of the alternatives rather than a feature of their composition. Whether the skeptic already uses the composition is where the argument must be careful, because it is too quick to say she does. A possible-worlds semanticist can hold that $A \vee B$ is true iff $A$ is true or $B$ is true and posit no worldly object composed of $A$ and $B$; using disjunction does not by itself commit one to composition on the world side. The composition is a genuine metaphysical commitment, and the paper asks for it openly. The claim is therefore conditional, and its interest lies in the conditional: \emph{once} a composition of worldly alternatives is admitted---as any reader who takes alternatives of reality to compose at all has admitted it---the relocation says that exclusivity need not be built into it. On such a reader's own view the composition she uses is the composition plus exclusions, and the relocation has told her where the exclusions live; decline them and the composition remains, now without the content, as a composition whose exclusion profile is now agreeing and not yet as the irreducible agreeing superposition, which the marker supplies. Two operations are distinct here: removing exclusion declarations from an already specified composition yields an agreeing exclusion profile, the composition and its alternatives retained throughout; specifying an unmarked filling, by contrast, does not yet specify that composition or its splitting regime. There is no further premise at which to balk provided the composition is granted, and the paper's thesis is exactly the conditional with that proviso foregrounded, not the claim that every classical theorist secretly held the antecedent. What silence does not do is \emph{produce} composition: absence of exclusion is not co-presence, and a reader who wanted co-presence to follow from silence alone would be right to refuse it. It follows from silence together with the composition the silence leaves unobstructed.

Second, the difference between the regimes is not notational, and this can be made exact:
The regimes are then defined on the structured objects, under the names the framework gives them, and the paper's central claim is a small conditional theorem about them rather than a slogan. Write $J(A_1, \ldots, A_n)$ for the primitive composition and $\mathrm{Ex}(A_i, A_j)$ iff $A_i$ excludes a pin $A_j$ holds, that is, iff $\mathrm{Int}(A_i)$ fills some attribute with a set omitting a pin that $\mathrm{Int}(A_j)$ holds; for a pair, $\mathrm{Agr}(A_i, A_j)$ iff neither excludes a pin of the other, and $\mathrm{Dis}(A_i, A_j)$ iff each does.

\begin{definition}[Agreement structure: agreeing and disagreeing superposition]\label{def:agreement}
The \emph{exclusion profile} of a primitive composition $J(A_1, \ldots, A_n)$ is the pattern of interior exclusions among its alternatives. The profile is \emph{agreeing} when $\mathrm{Agr}(A_i, A_j)$ for all $i \ne j$---no alternative's interior excludes a pin another's holds; \emph{disagreeing} when $\mathrm{Dis}(A_i, A_j)$ for all $i \ne j$---every pair excludes each other, each interior omitting a pin the other holds; and \emph{constrained} for every profile between. A composition is called an agreeing, disagreeing or constrained superposition \emph{in the profile sense} according to its profile; this is the sense classified below, and it is distinct from the irreducible agreeing superposition of Definition~\ref{def:irreducible-level}, a filling carrying the marker. In the vocabulary of regimes: the agreeing profile is the joint regime, the disagreeing profile the fully exclusive regime, and the constrained profiles the intermediate regimes---the regime names being the explanatory translations of the canonical terms, and not the reverse.
\end{definition}

\noindent ``Fully exclusive'' does not mean merely that the alternatives cannot be jointly asserted: a one-sided profile can also accompany incompatibility, as the marked example below demonstrates, and the exclusion profile alone does not determine whether joint assertion is invalid. It means that exclusion content is distributed symmetrically over every pair, and the distinction is one of update behavior: learning one alternative of a disagreeing profile eliminates every other, while a one-sided profile lets an already-holding alternative survive the learning. The regime names record where the exclusion content sits and how the composition updates, not whether joint assertion fails, which it does in both.

\begin{definition}[Three levels of the structure]\label{def:irreducible-level}
Three objects are distinct from here on. An \emph{unmarked filling} $[X : \{a_1, \ldots, a_n\}]$, the long paper's bare filling, specifies the available pins without a splitting regime and without a marker; the same filling with the attribute's carving is the \emph{exclusive splitting}; and the same filling with the irreducibility marker $e$, prohibiting narrower fillings, is the \emph{irreducible agreeing superposition}. An exclusion profile, of Definition~\ref{def:agreement}, is a property of a composition of alternatives with interiors; an unmarked filling has no such alternatives and so has no profile. With its carving it corresponds to the composition of its determinate cells, each a single-pin filling; the marked filling is not mapped to a composition here, and enters the results below only as an alternative within one, as in the worked example of the one-statement test.
\end{definition}
\noindent The three names are used with exactly these meanings throughout: \emph{agreeing profile} for the first level, \emph{unmarked filling} for the pins alone, and \emph{irreducible agreeing superposition} for the marked filling; where this paper writes ``agreeing superposition'' without qualification it means the profile sense, and the irreducible composition is always named as such. Agreement in the sense of Definition~\ref{def:agreement} is a property of a profile, the absence of interior exclusions, and an unmarked filling has neither an agreeing nor a disagreeing profile until a carving or a marker is intersected with it. The Regime Theorem classifies exclusion profiles, and so sorts compositions in the profile sense; it does not classify the marker, and the irreducible agreeing superposition is not among what it sorts. Declining exclusions therefore yields an agreeing profile and not the irreducible agreeing superposition, which needs the marker; the exclusivity thesis below, that exclusion is stated content, concerns the profile and is unaffected.

\noindent The one-statement test below exhibits that the constrained cases are well-formed objects that classical or-composition does not provide, although a structured encoding, a graph of exclusions for instance, can represent them. Two terms are deliberately absent from this paper: \emph{coherent} and \emph{decoherent}. Agreement structure is a logical category, defined on interiors before any weight, phase, or dynamics exists; coherence is a quantitative property of a physical realization of that structure, and the paper that installs the physics is where it is defined \parencite{AethusQuant}. No biconditional between the two vocabularies is asserted anywhere in the corpus, and the philosophical point of the present paper is that the category makes sense before quantum physics enters.

\begin{theorem}[Regime]\label{thm:regime}
Given the primitive composition $J$ and the interior profile $\{\mathrm{Int}(A_i)\}$, the regime of $J(A_1, \ldots, A_n)$ is completely determined by the profile: $(J, \{\mathrm{Int}(A_i)\}) \mapsto \text{regime}$. No second composition primitive distinguishes the joint from the disagreeing regime; changing only the interiors carries one and the same composition from joint, through the constrained cases, to fully exclusive.
\end{theorem}
\begin{proof}
The regime is defined by a condition on the interiors alone, quantifying over pairs; the composition operation does not enter the condition. So two compositions of alternatives with the same interior profile are in the same regime, and the same composition is moved between regimes by editing interiors and by nothing else.
\end{proof}

\noindent What the theorem consumes of $J$ is minimal and should be listed: that the composition exists; that it preserves and gives access to the interiors of its components; that adding a genuinely new statement to an interior changes the composed object; and that restating an already-present interior statement leaves it unchanged. Commutativity, associativity, and idempotence are inherited from the companion algebra as recorded in \S\ref{sec:relocation}, and the theorem does not consume them.

\begin{proposition}[The one-statement test]\label{prop:onestatement}
Joint and exclusive compositions, in the domain of Theorem~\ref{thm:regime}, differ in what a single added statement does. Add to the joint composition a single interior entry within $A_1$, the filling of $A_2$'s attribute that omits $A_2$'s pin---for a two-pin attribute, its complementary pin---and the composed object changes: its silence is broken, the pair's co-presence now constrained, with the constraint carried in $A_1$'s interior --- a coherent intermediate composition, note, that classical or-composition does not provide, exclusion there being all-or-nothing, though a structured encoding such as a graph of exclusions can represent it. The entry's semantics is the companion's complement mechanism \parencite{AethusTree}: stating in $A_1$ the filling that omits $A_2$'s pin refines $A_1$ by that filling, $A_1$ becoming $A_1 - A_1 A_2$, the part of $A_1$ compatible with $A_2$'s pin removed; $A_2$ is untouched. Where $A_2$ holds that pin alone in its attribute, as in the construction here, or carries the marker where it holds more, the two are now incompatible, their joint assertion vanishing; where $A_2$ permits narrowing, the omission removes only the overlap at the omitted pin. The test itself needs only that adding new interior content changes one structured object while restating existing content leaves another unchanged; the incompatibility claim is separable and holds under those hypotheses. Where it holds, incompatibility is symmetric; what the direction of the entry decides is which alternative gives up the overlap, and that shows in updating.

\noindent\emph{What the direction changes.} The profiles differ in what an update does, not in which alternatives are compatible, and the cleanest case has one attribute $X$ with two pins. Let $A_1 = [X : \{a_1\}]$, the state saying $a_1$ and not $a_2$, and $A_2 = [X : \{a_1, a_2\}] \cdot e$, the state holding $a_1$ and $a_2$ in irreducible agreeing superposition, the marker $e$ written explicitly, since the unmarked filling $[X : \{a_1, a_2\}]$ does not by itself prohibit narrowing and its incompatibility with $A_1$ would not be established without it; give them weights $w_1$ and $w_2$. $A_1$ excludes $a_2$, which $A_2$ holds, and $A_2$ excludes nothing $A_1$ holds: the one-sided profile. (The update that keeps every filling containing $a_1$ is learning that $a_1$ remains among the open possibilities; learning that the outcome is $a_1$ is a different observation, and the arithmetic here is the first.) The two-sided profile is $[X : \{a_1\}]$ with $[X : \{a_2\}]$, each omitting the other's pin; and the pins composed in one interior, $[X : \{a_1, a_2\}]$ alone, is the single filling of Eq.~\eqref{eq:forms}'s joint display, a different construction from a composition of alternatives. Learning means learning a pin's presence: to learn $a$ is to assert that $a$ is in the filling of its attribute, which keeps every alternative whose filling holds $a$ and drops every alternative whose filling omits it. The one-sided and two-sided profiles agree that their alternatives cannot be jointly asserted, in the two-sided case because $\{a_1\}$ and $\{a_2\}$ exclude each other, and in the one-sided case because $A_2$ carries the marker $e$ on $\{a_1, a_2\}$, which prohibits the narrower filling $\{a_1\}$, so that the joint-assertion rule returns invalidity; without the marker the fillings $\{a_1\}$ and $\{a_1, a_2\}$ would not contradict, which is why the marker is in $A_2$. The two profiles disagree on learning. Learning $a_1$ keeps both alternatives of the one-sided profile, $A_2$ standing at weight $w_2/(w_1 + w_2)$, and only $A_1$ of the two-sided; learning $a_2$ keeps only $A_2$ in both. So the direction shows up as an asymmetry of updating: the alternative that omits the other's pin is excluded by learning that pin, and the alternative that holds both pins is not. The profiles differ in content as well as in the direction of exclusion, $A_2$ holding $a_1$ in the one-sided profile and not in the two-sided one, so what the example shows is that two profiles with the same pairwise incompatibility support different updates, not that the direction of an exclusion alone, content held fixed, produces the difference. With weights $2$ and $3$, learning $a_1$ leaves $A_2$ the weight $3/5$ under the one-sided profile and $0$ under the two-sided; the one-sided profile matches the exclusive one on compatibility and on learning $a_2$, and the joint one on learning $a_1$, which is what constrained superposition is for a pair. The same pattern holds for alternatives over distinct attributes, an alternative stating the other's attribute with a filling that omits the other's pin. Add the same entry to the exclusive composition and nothing changes, the entry being already present. (Full pairwise exclusivity is the completion of such additions --- every interior written on --- and each step of the completion passes the same test.) A difference that one statement creates in one form and is idle in the other is a determinate structural difference within the structured-alternative model, not a difference of notation. In update notation, with $U_{i, \bar{a}_j}$ the operation adding to $\mathrm{Int}(A_i)$ the filling of $a_j$'s attribute that omits $a_j$: $U_{i,\bar{a}_j}(J_{\mathrm{joint}}) \ne J_{\mathrm{joint}}$, while $U_{i,\bar{a}_j}(J_{\mathrm{excl}}) = J_{\mathrm{excl}}$. What the test proves is nonidentity \emph{inside the representation}; whether the representation carves a mind-independent joint is the separate question \S\ref{sec:discussion} keeps separate.
\end{proposition}
\begin{proof}
In the joint composition the entry is absent, so its addition alters the composition's stated content --- the composed object before and after differ in what $A_1$'s interior says. In the exclusive composition the entry is present, so the addition is idempotent and the object unchanged. One operation, two outcomes, hence two objects.
\end{proof}
The test also fixes the direction of explanation. It is not that the joint composition is the exclusive composition with something mysterious subtracted; it is that the exclusive composition is the joint composition with something perfectly ordinary --- the exclusion of a pin --- \emph{added}, and added in a place (the interior) where additions demonstrably do work. The burden of the classical picture thereby inverts: on the relocation view, declining to state the negations does not leave composition undefined or degenerate --- it leaves the primitive composition in its agreeing regime in the bare sense, the absence of interior exclusion, which is permission to refine and not yet irreducibility, the open interior being the open door through which the composition already in hand passes unobstructed. What requires explanation in any particular domain is not why alternatives ever compose jointly, but why, so often, the interiors we find are already written on. (The companions have an answer --- the ubiquity of records: stated correlated structure whose joint invalidity with the alternatives eliminates their joint compositions, leaving the exclusive form as the residue, where a \emph{separating record} is a correlation whose branch states are orthogonal and which lies outside the reach of every admitted reversal, so that a partial correlation, with non-orthogonal branch states, and a reversible one, undone by an admitted operation, are not records and do not carve, not every physical correlation being an irreversible classical record, wholesale and for arbitrary alternatives, path-like or not. Exclusion is manufactured by elimination, not by a filling-in of values, and the classical disagreeing regime is the special case that ubiquitous record-keeping manufactures --- but that is their story \parencite{Aethus}.)

\section{What the Primitive Is, and Is Not}\label{sec:primitive}
Two boundary markers keep the object in focus, and the second is where the proposal's subtlety concentrates --- it is the fence that, once seen, prevents the two misreadings to which every classical habit will otherwise default.

First, the composition is a \emph{compositional} primitive---a primitive of the world-side algebra of alternatives, and not a new connective or a revision of logic, which the paper leaves untouched---and it is primitive in the same explanatory sense the logical constants are: usable, not further describable. One deploys \emph{and} and \emph{or} constantly without ever articulating their substance; there is nothing conjunction is \emph{like}, and the demand for such an account is recognized as a category mistake the moment it is made explicit. The agreeing regime asks for exactly the same treatment. To the question ``but what \emph{is} the jointness of jointly composed alternatives?" the correct answer is the one we already accept for conjunction: it is what the composition \emph{does}, exhaustively specified by its role in the algebra of Eq.~\eqref{eq:forms} and the test of Proposition~\ref{prop:onestatement}, with no residue awaiting phenomenological description. That the primitive feels intangible is not a defect of the proposal; intangibility of exactly this kind is the signature of the logical constants, and the claim that our inventory of them was complete was only ever an induction from the classical cases.

Second --- the load-bearing fence --- whatever the occupants of a joint composition may \emph{do} in any given domain is no property of the jointness itself. The primitive supplies \emph{openness}: it ushers the allowance of co-presence, opening a compositional void where the exclusive form would have walled each alternative into solitude. What, if anything, \emph{uses} that openness --- what fills the opened room, what phenomena the co-present alternatives jointly sustain --- is a further fact about the occupants and their domain, strictly downstream of the regime and never to be read back into it. Jointness is an abstract, attribute-mediated state; use is material, domain-relative, and separately theorized.

The fence blocks two misreadings, opposite in direction and equally natural. The first inflates: it takes jointness for a kind of causal glue or interaction-substance, a medium \emph{through which} co-present alternatives influence one another --- and then rightly complains that no such medium has been characterized. But the complaint attacks an addition the proposal never made: the regime licenses co-presence and nothing further, and any influence its occupants exhibit belongs to their own story, not to the license. The second misreading deflates: it takes interior silence for mere ignorance --- the alternatives are joint ``for us," because we have not yet learned which one obtains. This mistakes an asserted structural fact for an epistemic gap. The silence of Eq.~\eqref{eq:forms}'s left form is \emph{part of the composition's content}, as the one-statement test shows: writing in the negation \emph{changes the composed object}, which distinguishes the structured representations, though epistemic representations also change when information is added, so the test establishes distinguishable representations and not, by itself, their worldly interpretation. Interior silence is constructive, not confessional --- a claim about what the composition \emph{is}, not about what we know of it.

That point deserves its own record, since it is where the proposal touches, and departs from, its nearest neighbors:
\begin{remark}[Silence as compositional determinacy]\label{prop:silence}
Interior silence is not indeterminacy \emph{of} the alternatives but determinacy \emph{about} their composition: a joint composition determinately composes jointly, and the absence of interior negations is an exact structural fact about it, fixed by Proposition~\ref{prop:onestatement}'s test, not a deficiency awaiting resolution.
\end{remark}

\section{The Sister Taxon}\label{sec:taxon}
The category claim can now be stated at full strength. The agreeing regime is not a variant, limiting case, or exotic corner of anything classical metaphysics ever catalogued. The familiar physical \emph{or} --- alternatives excluding by nature, exactly one obtaining --- is, on the relocation view, entirely a subcategory of the \emph{exclusive} form: it is what composition looks like when the interiors are fully written on. The joint regime is not a member of that family at all. The right image for the discovery-shape is taxonomic: not a new species within a known genus, but a sister taxon to the entire classical catalogue --- a second fundamental way for alternatives to stand to one another, coordinate with, rather than subordinate to, the way we knew. The algebraist's vocabulary makes the taxon's signature exact, (the terms carrying their standard meanings in the companion's weighted composition algebra \parencite{Aethus}): a fully resolved alternative is \emph{simple} --- maximally structureless --- while the irreducible agreeing superposition, the filling carrying the marker of Definition~\ref{def:irreducible-level}, is \emph{indecomposable without being simple}: admitting no decomposition into exclusive parts, a decomposition here meaning a splitting into narrower fillings, which the marker prohibits, yet carrying arbitrarily much internal structure in its stacking. That description belongs to the marked construction and is not a consequence of Theorem~\ref{thm:regime}: an agreeing profile alone supplies no irreducibility, the marker does. Nothing classical occupies that cell of the classification, which is the taxonomic claim in one line.

One neighbor must be distinguished before the taxon claim can stand, because it too countenances a disjunction true with neither disjunct true: quantum logic \parencite{birkhoffvonneumann1936, putnam1968}. The difference is where the revision lands. Quantum logic obtains its non-classical disjunction by \emph{revising the connectives} --- surrendering distributivity, so that the lattice of propositions itself is non-Boolean. Here no connective is touched and the logic stays classical throughout: what varies is what the alternatives \emph{state} --- whether the interiors carry negations --- so the non-classicality lives in noncontradiction's \emph{domain}, never in its laws. The agreeing regime is not a rival logic; it is a second way for classically-logical alternatives to compose, and the contrast with Birkhoff--von Neumann is thereby a stated one rather than an omission.
% >>> EDIT U1 (Fable 5, Aug 29 2026): the differentia restated at the grade a quantum logician will insist on
Stated at the grade a quantum logician will press for, the differentia is \emph{where} the revision lands and not whether there is one: the present proposal revises the primitive stock of composition while leaving the connectives and the propositional lattice alone, and a proposal that adds a composition primitive is a revision by anyone's accounting. Putnam's question accordingly does land here, and receives a changed answer --- the empirical matter is the primitive inventory rather than the consequence relation, and, downstream of it, which exclusions the alternatives state.
% <<< END EDIT U1
The nearest existing family in the literature deserves its coordinates, both as credit and as contrast. The determinable-based account of metaphysical indeterminacy \parencite{darby2010, skow2010, wilson2013, calosi2019} countenances states in which a system possesses a determinable property while lacking any determinate value of it --- a genuine structural cousin of an interior-silent alternative set, and the closest thing in print to the present proposal's states. The contrast is not that the present proposal supplies a composition story where the literature has none; that claim would be false. The situation-semantic model of Darby and Pickup builds indeterminacy from incomplete, local situations and their \emph{fusion}, with explicit restrictions on composition, and applies it to superposition and entanglement, its central question being how truth persists through such fusions without contradiction \parencite{darbypickup2021}; and the determinable-based program has been developed for double-slit interference specifically, with a glutty implementation on the table \parencite{calosiwilson2021}. Composition, fusion, precisification, and determinable structure are all in print. The differentia is in the \emph{kind} of composition story. Those accounts treat the composite primarily through indeterminacy---gaps, gluts, restricted persistence of truth across the fusion---and the composite is where the revisionary work happens. The present proposal makes the composition itself \emph{determinate}, a determinate object of a non-classical kind, and relocates the revisionary content into the alternatives' interiors as removable exclusion. Darby and Pickup are the nearest neighbours by some distance: partial situations standing beside one another, fused, with quantum applications and disagreement without contradiction. The difference between their fusion-with-restricted-persistence and the present joint-composition-with-interior-silence is exactly this---for them what is restricted is how much of a situation's content survives fusion, while here nothing is restricted and nothing survives partially, the interiors being either written or silent and the composite determinate either way. Eq.~\eqref{eq:forms} and Proposition~\ref{prop:onestatement} are the chapter that difference needs: the two composition forms, and the one-statement boundary between them. And Remark~\ref{prop:silence} marks the departure: on the present view the joint composition is not \emph{indeterminate} in any sense that calls for a revisionary logic or a gappy metaphysics \parencite{barneswilliams2011} --- it is a determinate composition of a non-classical kind, and the appearance of indeterminacy is an artifact of insisting on describing the sister taxon in the home taxon's vocabulary.

The engagement owes one sharpening at its most exposed point. Calosi and Wilson run the determinable strategy on \emph{entanglement} itself \cite{calosi2019}: the entangled subsystem determinately has the determinable while lacking any determinate value --- indeterminacy located as a \emph{gap} in the determinate row. That is composition-adjacent, and it is the sentence against which the differentia must be stated; so it is stated at the same address. The determinable account gaps the \emph{value} and leaves exclusion's domain untouched: the alternatives remain rivals, one determinate owed, the owing merely unmet. The present relocation moves \emph{exclusion itself} --- interior-stated, or-jointness primitive, co-presence with openness --- so that in the agreeing regime nothing is owed and nothing is gapped: there is no missing determinate, because the regime never contracted for one. Gap versus relocation is the whole difference, and it is a difference in \emph{where a law lives}, not in how much indeterminacy to admit.

One clarification protects the claim's precise reading. Nothing here disputes that individual open states have been described before --- the determinable literature does exactly that, and does it well. The claim is about \emph{composition as such}: that co-present alternatives admit a joint mode of standing together, governed by its own side of Eq.~\eqref{eq:forms}, has to the author's knowledge never been articulated as a category within a realist metaphysics, and it is that compositional thesis, not the bare possibility of openness, that constitutes the sister taxon.

\section{The Color-Exclusion Problem, Reclassified}\label{sec:color}
Run the machinery on the founding case. The datum: \emph{red-at-p-at-t} and \emph{green-at-p-at-t} exclude. The classification problem: the exclusion is not logical (no shared negated conjunct; Wittgenstein's 1929 concession stands) and not comfortably empirical (its necessity-feel survives every imagined experiment). The relocation view supplies the third classification the datum was demanding: the color-alternatives \emph{as we in fact find them} carry the mutual negations as interior content --- the exclusion is \emph{stated structure}, contingent in the sense that the interiors could have been silent, structural in the sense that, being written, it is part of what the found alternatives \emph{are}. Contingent-yet-structural is precisely the profile that resisted the old dichotomy: too necessary-feeling for brute empirical generality, because it is constitutive of the alternatives as constituted; too removable for logical form, because constitution by stated content is constitution that could have been otherwise.

The diagnosis of the 1929 failure follows in one sentence. Wittgenstein tried to push exclusion \emph{into logical form} --- to make the connectives carry it --- and the attempt failed because exclusion is not the connectives' work; the repair the problem actually wanted pushes exclusion into \emph{stated content}, which is where a removable thing must live. The Tractarian instinct that exclusion had to be \emph{somewhere structural} was correct; the error was an impoverished inventory of structural locations, containing logical form and the world's brute furniture and nothing between. The interior of an alternative is the between.

The conditional reclassification also pays a modal dividend, and it is a mark in the proposal's favor rather than a bullet bitten. On the relocation view, joint color-occupation is \emph{conceivable} --- the interiors could have been silent, and Eq.~\eqref{eq:forms}'s left form describes the result without strain --- while its absence from the found world is a fact about which compositions the world's alternatives happen to instantiate, not about what composition permits. This gets the phenomenology exactly right. The color-exclusion datum was always double: the exclusion feels necessary, \emph{and} the excluded scenario feels imaginable --- painters, dreamers, and philosophers keep almost-picturing it. The classical classifications must convict one half of the datum of confusion; the relocation view convicts neither. The exclusion is real and structural where it obtains; the imagined alternative is a well-formed composition the found interiors decline to instantiate. A scope line keeps the modal claim honest, and it is sharper than a concession to essentialism. What the declension argument secures is the coherence of the joint \emph{form}: a pair of positive qualities structurally like red and green, with silent interiors, would compose jointly. Whether \emph{these} determinates---red and green as constituted---could have done so is a further claim, and it turns on a question the paper has not answered: property identity across changes of interior content. Write actual red as $\langle \text{positive red character},\ \neg G, \neg B, \ldots \rangle$ and remove the exclusions to obtain $\langle \text{positive red character} \rangle$; whether the second is still \emph{red}, rather than a distinct property resembling it, is not settled by anything above. If the exclusionary interior is essential to red, actual red cannot enter the agreeing regime and the instance-level modal claim fails; if it is accidental, criteria are owed for what individuates the property independently of that interior. Two theories of individuation are available, and the paper's results divide across them. On \emph{core individuation}, $A = B$ iff $\mathrm{ch}(A) = \mathrm{ch}(B)$, so that changes of exclusion profile preserve property identity and $\langle \text{red character} \rangle$ is still red; on \emph{full-structure individuation}, $A = B$ iff $\mathrm{ch}(A) = \mathrm{ch}(B)$ and $\mathrm{Int}(A) = \mathrm{Int}(B)$, so that removing an exclusion makes a different property. The Regime Theorem, whose one-statement test is a criterion of structured-state identity rather than of property identity, needs neither; the modal claim that red itself could coexist with green needs the first, and is therefore stated as a \emph{Color Corollary conditional on core individuation}: under full-structure individuation the general composition thesis stands and the red--green corollary fails. The paper takes no stand between them, because the question is the same as its booked open problem---where an alternative's interior content ends---and that problem is more load-bearing than a booked item usually is: until the boundary of an alternative's content is principled, ``the exclusion is interior to red'' is itself underdetermined. So the section's title is \emph{reclassified} and not \emph{dissolved}: the datum receives the third classification it was demanding, contingent-yet-structural, and the traditional problem at full strength---red and green \emph{themselves}---waits on the individuation question. Section~\ref{sec:exposure} assigns the instance-level question to the witness rather than to the armchair. One feature of the found writing deserves its own sentence: color exclusion is \emph{systematic} --- red excludes every determinate of the family at once, not green alone --- and this is a datum about \emph{how} the interiors are written, not an embarrassment to the classification. Stated content can be stated wholesale: the determinable structure of the family \parencite{wilson2013} is exactly the packaging by which an interior carries a whole family's mutual negations in one stroke, so the systematicity of the found exclusions reflects the systematicity of their inscription. Systematic is not the same as necessary, and the relocation view needs only the former.

\section{The Direction of Discovery}\label{sec:exposure}
One fact about this proposal's history belongs in the paper, because it bears on provenance, a kind of record philosophical theses rarely have. The two composition forms of Eq.~\eqref{eq:forms} were not first produced by philosophical argument and then sought in the world. They arose first from the author's analysis of physical phenomena within the developing framework, before their present philosophical application. In the companion framework \parencite{Aethus, AethusBrief}, a classification theorem proved over the powerset of co-present alternatives, from that framework's own postulates, generates exactly this paper's dichotomy: its \emph{valid} class is the joint composition, its consistency \emph{pruning} is the manufacture of the exclusive composition, and within that framework's reading of interference phenomena the agreeing regime is the load-bearing cell.

 The companion framework is a proposal, and whether the world instantiates the agreeing regime is precisely what it proposes; nothing here launders that proposal into an observation.

 In the case-by-case analysis it is explicitly the \emph{interior} of each configuration that forces the nonexistence of certain traversal instances and thereby --- strictly indirectly --- blocks whole elements of the powerset from realization: exclusion visibly operating as interior content, doing its work by statement rather than by nature. The relocation defended here was that observation's \emph{universalization}, arrived at by pressing the inductive question in its original form --- \emph{what must a universe look like, inclusive of all peripheral context, which insists on providing this ontology upon the powerset?} --- whose answer is that such a universe is presenting the projection of a deeper principle about how classical exclusivity emerges as a special case.

\textit{The deductive core.} The goal was a coherent account of superposition, and that goal forced exclusivity to be \emph{removable}: agreeing superpositions in the profile sense are precisely co-present alternatives that do not exclude, the composition itself being the granted primitive. What is removable can be neither logical form, (logic does not come off), nor the alternatives' brute nature, (natures do not come off). Other locations remain in principle---an exclusion relation between the alternatives, as truthmaker semantics has it \parencite{fine2017truthmaker}; a law governing their co-instantiation; determinable structure; a constraint on admissible composites---and the argument does not eliminate them. What it claims is abductive: \emph{stated content} --- something the alternatives carry, written on their interiors, and therefore capable of going unwritten --- is the location that makes exclusion removable and leaves ordinary logic untouched at once, which no rival location does as economically, an exclusion relation being as much a primitive as the exclusivity it replaces. The relocation is not the sole survivor of an elimination; it is the location the analysis found most economical, and the sections above are the case that the economy is real.\footnote{The recognition was gradual rather than punctual, consolidating around the turn of 2022 to 2023 and established by May 2023, through repeated analysis of agreeing-superposition cascades in the double-slit case; the inferential infrastructure was assembled to isolate the principle's substance from the case that kept exhibiting it, and the interior-forcing observation above is that infrastructure's keystone. The companion record documents the development \parencite{Aethus}. The companion record also preserves the author's contemporaneous assessment of the idea's significance.} What the documented order establishes is therefore not confirmation but \emph{provenance}: the dichotomy arose from the author's analysis of the case-by-case structure of the phenomena within the developing framework, not invented for the philosophy and then sought in the world --- the philosophy was reverse-engineered from the forms, never the forms from the philosophy.

This bears on a question a careful reader will have been holding since Section~\ref{sec:relocation}, and bears on it at a stated strength: whether the relocation is a discovery about composition or an elegant piece of bookkeeping a metaphysician may adopt or decline. Within the paper's own resources the question splits: that the agreeing regime in its bare form is non-optional \emph{given} the relocation is the declension argument's work, the irreducible form needing the marker the argument does not supply; whether exclusion \emph{is} interior content, rather than merely representable as such, is a question provenance bears on without settling. And provenance of that kind is what the conventionality worry must reckon with, though it does not settle it: a distinction gerrymandered for a philosophical purpose has its shape explained by the purpose, while a distinction extracted from independent structure --- structure analyzed for entirely different ends, before the philosophical thesis existed --- has its shape explained by what it was extracted from. What that establishes is exactly this: the documented provenance reduces the probability that the distinction was tailored after the fact to the present philosophical application, and it establishes neither that the companion framework is correct nor that the distinction carves at a joint rather than being a representational choice. The quarantine of Section~\ref{sec:discussion} stands exactly as stated --- no thesis of Sections~\ref{sec:relocation}--\ref{sec:color} takes the witness as a premise, and a reader who rejects the companion physics wholesale loses nothing of the paper's arguments. What such a reader loses is only the answer to \emph{this} section's question: with the witness, the relocation is what the framework's analysis of the phenomena arrived at before any philosophical application; without it, the relocation is an elegant hypothesis awaiting one.

\section{Commitments, Objections, and Placement}\label{sec:discussion}
\emph{What is claimed.} Four theses, in descending order of strength: that exclusion among co-present alternatives is stated interior content rather than a nature of disjunction (the relocation); that the agreeing regime falls out of the relocation with no further postulation, and is fixed against the disagreeing regime by a one-statement test (the declension and the test); that the resulting primitive supplies openness only, with every use-fact quarantined to the occupants (the fence); and that this reclassifies the color-exclusion datum and supplies the determinable literature with composition (the applications).

\noindent Four claims travel under the paper's title, and they have four different statuses; the table fixes them so that the formal result is not credited with the others.
\begin{center}
\begin{tabular}{p{0.66\linewidth} p{0.26\linewidth}}
\toprule
\textbf{Claim} & \textbf{Status} \\
\midrule
\textbf{Formal}: an interior-exclusion representation distinguishes joint, constrained, and disagreeing regimes of one composition & proved (Theorem~\ref{thm:regime}, Proposition~\ref{prop:onestatement}) \\
\textbf{Metaphysical}: actual exclusion is in fact interior structure & argued here \\
\textbf{Instantiation}: primitive joint composition exists in nature & proposed by the companion physics \\
\textbf{Identity}: removing exclusions can preserve the identity of familiar properties & required by the color corollary; conditional on core individuation \\
\bottomrule
\end{tabular}
\end{center}


\emph{What is not claimed.} That the agreeing regime is physically instantiated. That thesis belongs to the companion works \parencite{Aethus, AethusBrief}, where the joint composition of Eq.~\eqref{eq:forms} is put to load-bearing use;\footnote{One interplay deserves its coordinates even under the quarantine, because it is close to the bone. The companions prove a classification theorem over the powerset of co-present alternatives, and the two results are nearly one: the theorem's \emph{valid} elements are exactly the joint compositions of Eq.~\eqref{eq:forms}'s left form, and its consistency \emph{pruning} is exactly the machinery by which the exclusive form gets manufactured --- the interiors doing the blocking, and doing it, notably, in a strictly \emph{indirect} sense, each interior forcing absences that only thereby foreclose other whole elements. Historically the dependence ran the other way: the case-by-case powerset analysis was the \emph{inductive root} from which the present relocation was first universalized --- a universe that supplied the powerset classification in isolation would, by an Occam argument, be presenting the projection of a deeper principle about how classical exclusivity emerges as a special case, and the relocation is that principle stated. The present paper is the philosophical face of the same structure, argued without the theorem; the companion is the combinatorial face, proved without the philosophy. (A third face completes the program: a perturbation semantics for counterfactuals built on the same states and the same pruning, the three surveyed together in \parencite{AethusBrief}.)} no argument of the present paper leans on it, and the paper's theses stand or fall on the philosophical ground above --- with Section~\ref{sec:exposure}'s witness entering only at the level of the theses' \emph{status}, never as a premise of any of them. One sentence of connection is owed and no more: \emph{if} the regime is instantiated, then any interference-grade phenomena its occupants exhibit would be, by Section~\ref{sec:primitive}'s fence, the occupants' \emph{use} of the openness --- never the openness itself, and never evidence about the primitive's nature.

\emph{One placement the companions have since made, recorded here so the two composition doctrines do not appear to compete.} The primitive of this paper governs how \emph{alternatives} compose---the jointness of \emph{or}. A second commitment about composition, named in the companion works on August 29, 2026 as \emph{Decomposition Neutrality} and traced to the framework's 2021 Attribute Set, governs \emph{conjunctive} composites instead: a stated-attribute has no privileged decomposition, so a conjunction's answer space is the product of its parts' and never a coarsening of it, which is the structure on which the counterfactual semantics reads a conjunctive antecedent's reopening, though its dual-pin reading there is a substantive interpretation of the antecedent that Decomposition Neutrality does not by itself supply \parencite{AethusBrief, Aethus}. The two meet nowhere: one is a posit about how alternatives are held together, the other a constraint on how a conjunctive attribute's states are counted, and the companions' ledger books them as separate entries---the union as the declared inductive fulcrum, the constraint at the grade of the record-individuation principle it sits beside.

\emph{Objections, pre-empted.} \emph{Negative facts through the back door}: the interior negations of the exclusive form are stated attribute-content --- things the alternatives \emph{say} --- not worldly negative facts; the proposal's ontology of negation is exactly as innocent as its opponent's ontology of exclusion, since the classical view also needs the world to somehow enforce $\neg(A_1 \wedge A_2)$, and locates the enforcement less perspicuously. \emph{Mere notation}: answered by Proposition~\ref{prop:onestatement}, with its status stated exactly --- the test is a formal fact about the proposal's own objects, establishing that the two forms differ structurally \emph{within} the framework; a thoroughgoing deflationist who rejects the objects wholesale is not refuted by it, only denied the cheaper complaint that the difference was verbal. \emph{Why call it ``or" at all}: because the joint form is what remains of disjunction when its classical costume is removed --- the union of Eq.~\eqref{eq:forms} \emph{is} the inclusive disjunction Jevons left us, and the naming follows the algebra; the reader who prefers a new word loses nothing but the historical continuity, which is real. \emph{Non-contradiction}: fully preserved, and the mechanism deserves statement since it completes the logical-realist bookkeeping --- what remains impossible on the relocation view is the configuration in which one alternative's content and its rival's are stated \emph{of the same subject in the same respect} --- the \emph{intersection} of disjoint alternatives; jointness is union, never intersection, and the classical intuition of impossibility attaches, rightly, to the latter.

\emph{The open problem, booked --- and scoped exactly.} Within the formalism nothing here is open: \emph{interior} is a scope notion and nothing more --- content is interior to $A_1$ just when it sits within $A_1$'s own term of the composition, inside what the outer union joins --- so the displays of Eq.~\eqref{eq:forms}, the one-statement test, and every argument above are fully determinate as they stand. What is booked is the \emph{application} question: when the alternatives are worldly rather than schematic --- \emph{red}, \emph{green} --- what fixes which found content belongs to an alternative's stated interior rather than merely accompanying it, i.e., where one alternative's content ends. That is an instance of the general problem of individuating content, not a defect peculiar to the relocation, and the companion framework's record-based account of when a distinction exists to be drawn is the natural place its answer would come from \parencite{Aethus}. It is booked as open here rather than answered, in the judgment that an honest booking is worth more than a definitional discharge.

\emph{Placement.} Against the enriched-language tradition: the content is agreed, its location moved from semantics to the alternatives themselves. Against necessitarianism about exclusion: structure is agreed, ineliminability denied, and the one-statement test supplies what necessity never could --- an account of what the difference between excluding and not excluding \emph{consists in}. Against the determinable framework: cousinhood acknowledged, composition supplied, indeterminacy reclassified as non-classical determinacy. And against the Tractarian episode itself: the instinct honored, the inventory enlarged, the problem not solved but conditionally \emph{reclassified}, so that its stubbornness is explained rather than outlasted, on the property-identity qualification stated above. Fine's truthmaker framework holds the nearest \emph{primitive} rival on the question of exclusion's address \cite{fine2017truthmaker}: exact exclusion as a basic relation between states, negation built upon it. The relocation thesis answers the same question derivationally: exclusion is never basic here, arriving only as stated interior negation --- content, on the books, removable by silence --- which is precisely what lets the agreeing regime exist at all, in its bare form, with irreducibility a further stated entry; a primitive exclusion relation, on which Fine's framework places conditions rather than declaring universal pairwise exclusion, refers the regime the companion classification finds occupied to those conditions, and whether they exclude it is an incompatibility still to be proved, while a stated exclusion prices it. And the closest \emph{formal} cousin deserves its name: possibility semantics \cite{humberstone1981from, holliday2021possibility}, models built from partial possibilities refined rather than completed --- partiality without forced totality, the shape this paper needs. The differentiae are three and exact: possibilities persist under refinement, where agreeing composition is \emph{not} a refinement of either component; the possibilities carry no evaluator index; and nothing weighs them. Partiality alone reaches the door; the index and the weights are what walk through it.

\paragraph*{Provenance note.} The framework this paper extracts from is the author's, its development documented in successive timestamped deposits \parencite{Aethus, AethusBrief}. The relocation was recognized gradually, consolidating around the turn of 2022 to 2023 and established by May 2023; the case-by-case analysis from which it was universalized is described in Section~\ref{sec:exposure} and documented in the companion record. AI assistance was used in preparing this manuscript to conventional presentational standards, in a recognition-and-translation capacity; the relocation, the primitive, and the theses defended are the author's, with their development dated in the companion corpus.

\paragraph{Fillings, carvings, and the two sums.} The pins inside a filling are the open possibilities for the attribute and say nothing about exclusivity: $[X : \{a_1, \ldots, a_n\}]$, an unmarked filling, has neither an agreeing nor a disagreeing profile until something is intersected with it. Intersected with the single-pin carving $c_X$, it is the union of the determinate states $[X : \{a_i\}]$, the canonical form of exclusivity with which this paper began, its outer union composing states; intersected with the irreducibility marker, its determinate children are tagged and it is the agreeing superposition proper. A carving's cells are the determinate proper children of the carved state, the terms of its decomposition $A = \sum_i w_i B_i$ whose weights are not measure-zero relative to the rest, which is what licenses the decomposition; adjoining a determinate assertion multiplicatively, forming $A \cdot [X : \{a\}]$ without the carving, yields a child in the multiplicative sense and declares no split. The two operations remain distinct, and the support semiring keeps them apart: assembling pins into a filling yields the one Aethus $[X : \{a, b\}]$, while collecting Aethae into a composition, the support union, yields $\{[X : \{a\}], [X : \{b\}]\}$, and $\{[X : \{a\}]\} \cup \{[X : \{b\}]\} \ne \{[X : \{a, b\}]\}$. What connects them is the carving, a map from the carved filling to the collection of its determinate states, not an equality, which the algebraic companion's expansion discussion denies in its contracted semiring; the three roles above are what the canonical form left open, now settled by naming the map.


\section*{Provenance}
\addcontentsline{toc}{section}{Provenance}

\noindent\textit{On Fillings, carvings, and the two sums.} Author's note (2026-09-25; quoted verbatim, its \emph{peg} being the present \emph{pin}): ``saying nothing about $b$'' and ``$b$ is blank with all pegs filled'' are two names for the same thing, such that ``$b$ is an ordinary blank'' or ``$b$ is an irreducible blank'' actually say nothing about $b$ at all---the irreducibility comes from the $b^{\ast}$ thing. [\ldots] We ought to think of pegs filling a state as simply saying ``these are the open possibilities for this value'', thereby being passive in terms of stating irreducibility or not, and instead simply putting out on the table what the superposition actually is. Then, it is by intersecting $b^{\ast}$ that we attain the Aethus for which the net concoction is that proper child Aethae making $b$ determinate are tagged invalid by $b^{\ast}$, or alternatively by intersecting that the carving has to be to determinate states that we can write the Aethus as equivalent to the union of the two determinate states. The blank itself is more similar to weighing options, even though it still is not real scenarios, because it admits no validity to any proper child Aethae in which $b$ is determinate anyway.

\printbibliography
\end{document}